\documentclass{leveldoc}
\DocUses{foundations,first_bend,slingshot,stray_field}
\DocFingerprint{971577a791eee5de}

\begin{document}

\section{What the level asks}\label{sec:asks}

The game's text:
\leveldescription

One ion is launched from $(\val{xZero}, \val{yZero})$ along $+x$ with $T_0 = \val{Tzero}$ in a
world with $c = \val{c}$, and must enter the detector
$\val{xDetLo} \le x \le \val{xDetHi}$, $\val{yDetLo} \le y \le \val{yDetHi}$. The player places up
to three antennas: oscillating electric dipoles $\vect{p}(t) = \vect{p}_0\cos\omega t$, all driven
by the level's generator at $\omega = \val{omega}$ and in phase, each with an amplitude from the
listed magnitudes, either sign (a negative one is the opposite phase), and one of the
orientations $0^\circ$, $45^\circ$, $90^\circ$, $135^\circ$ (\physref{2.4}). A flight ends when the
ion enters the detector, touches an antenna's body (a sphere of radius $\val{Ra}$), leaves the
world, or reaches $t = \val{tMax}$ (\xsecref{foundations}{sec:board}).

The reference places one antenna of amplitude $p_0 = \val{pA}$ at $(\val{xA}, \val{yA})$, oriented
at $\val{angleA}^\circ$ (along $+y$), directly below the launch. At $t = 0$ its dipole points up,
and the field above a dipole along its axis points along it: the ion is pushed up while it is
near, before the field reverses, and coasts into the detector. The document derives the
antenna's field from Maxwell's equations, marches the flight numerically, and explains the kick
to first order in the field.

\section{Setup}\label{sec:setup}

The level file, field by field (the appendix reproduces it):
\setuptable

The ion has charge $q = \val{q}$ and mass $m = \val{m}$. By \xeqref{foundations}{eq:launch} it
starts with $\gamma_0 = \val{gammaZero}$, momentum $p_\mathrm{i} = \val{pZero}$ and speed
$v_0 = \val{vZero}$. The drive's period is $2\pi/\omega = \val{period}$ and its wavelength
$2\pi c/\omega = \val{wavelength}$, longer than the board: the ion flies in the antenna's near
and induction zones ($kr = \omega r/c = \val{kr}$ at the launch, $\val{rLaunch}$ from the
antenna). The world is $-1 \le x \le \val{xWorld}$, $-1 \le y \le \val{yWorld}$.

\begin{figure}[htbp]
\centering
\begin{tikzpicture}
\begin{axis}[axis equal image, width=0.92\linewidth, xmin=-1.6, xmax=31.6, ymin=-1.6, ymax=21.6,
  xtick={0,5,...,30}, ytick={0,5,...,20}, tick label style={font=\small},
  xlabel={$x$}, ylabel={$y$}, label style={font=\small}]
  \draw[black!8] (0,0) grid (30,20);
  \draw[black!45, dashed] (-1,-1) rectangle (\rawval{xWorld},\rawval{yWorld});
  \fill[green!50!black, fill opacity=0.18] (\rawval{xDetLo},\rawval{yDetLo}) rectangle (\rawval{xDetHi},\rawval{yDetHi});
  \draw[green!40!black] (\rawval{xDetLo},\rawval{yDetLo}) rectangle (\rawval{xDetHi},\rawval{yDetHi});
  \addplot[thick, blue!70!black] table {pathref.dat};
  \addplot[thick, orange!80!black] table {pathanti.dat};
  \addplot[thick, teal] table {pathlate.dat};
  \addplot[thick, black!50, dashed] table {pathsame.dat};
  \fill[red!75!black] (\rawval{xA},\rawval{yA}) circle[radius=\rawval{Ra}];
  \draw[->, thick, red!75!black] (\rawval{xA},\rawval{yA}) -- ++(axis direction cs:0,1.8);
  \fill[orange!80!black] (\rawval{xAnti},\rawval{yAnti}) circle[radius=\rawval{Ra}];
  \fill[teal] (\rawval{xLate},\rawval{yLate}) circle[radius=\rawval{Ra}];
  \fill (\rawval{xZero},\rawval{yZero}) circle[radius=0.15];
\end{axis}
\end{tikzpicture}
\caption{The reference antenna (red, its dipole at $t = 0$ drawn) and its flight (blue). The other
flights of \S\ref{sec:others}: an antenna at $(\val{xAnti}, \val{yAnti})$ in the opposite phase
(orange) and in phase (grey, dashed: it leaves the world), and one at $(\val{xLate}, \val{yLate})$
in phase (teal).}
\label{fig:board}
\end{figure}

\section{Principles and model}\label{sec:model}

\begin{principles}
  \item Maxwell's equations \xeqref{foundations}{eq:maxwell} and their retarded solution in the
    Lorenz gauge: the potentials
    $\phi(\vect{x},t) = \int \rho(\vect{x}', t_\mathrm{r})\,\dd^3x'/R$ and
    $\vect{A}(\vect{x},t) = c^{-2}\int \vect{J}(\vect{x}', t_\mathrm{r})\,\dd^3x'/R$ with
    $R = \abs{\vect{x}-\vect{x}'}$ and $t_\mathrm{r} = t - R/c$ (\Griffiths{Ch.~10},
    \Jackson{Ch.~6}, in the units of \xsecref{foundations}{sec:units}), and the fields
    $\vect{E} = -\nabla\phi - \partial\vect{A}/\partial t$, $\vect{B} = \nabla\times\vect{A}$.
  \item The fields of an oscillating electric dipole (\Jackson{§9.2}, \Griffiths{Ch.~11}),
    derived below in the time domain.
  \item The relativistic equations of motion \xeqref{foundations}{eq:eom} with the launch
    \xeqref{foundations}{eq:launch}, and Liénard's radiated power \xeqref{foundations}{eq:larmor}.
  \item The board's rules; the classical Runge--Kutta method, cubic Hermite interpolation and
    the error estimate from the step $2h$ (\xsecref{foundations}{sec:board},
    \xsecref{foundations}{sec:numerics}), as in \xsecref{slingshot}{sec:method}.
\end{principles}

The antenna is a point dipole at $\vect{r}_\mathrm{a}$ with
$\vect{p}(t) = \vect{p}_0\cos\omega t$ for all $t$, past and future: the generator has been
running long before the launch, so the field at the launch is already the steady oscillation.
The ion moves by \xeqref{foundations}{eq:eom} in this field, from
$\vect{x}(0) = (\val{xZero}, \val{yZero})$, $\vect{p}(0) = (p_\mathrm{i}, 0)$, until the first of
the board's events.

\begin{assumption}{The ion's own radiation}\label{as:radiation}
The ion is accelerated and radiates the power \xeqref{foundations}{eq:larmor}; the game neglects
the energy this carries away (the level has no radiation reaction), and so does this document.
Along the reference flight it is $\val{radRef}$, a fraction $\val{radFrac}$ of $T_0$
(\S\ref{sec:work}).
\end{assumption}

\begin{assumption}{The antenna is driven}\label{as:driven}
The dipole oscillates with a fixed amplitude whatever the ion does: the generator supplies the
power it radiates, $p_0^2\omega^4/(3c^3)$ (\physref{2.4}), and the ion's field does not act back
on it. This is the game's model; a real antenna's current is set by its generator and its
impedance, and a passing ion's influence on it is negligible.
\end{assumption}

\section{Derivation}\label{sec:derivation}

\subsection{The antenna's field}\label{sec:field}

Put the dipole at the origin, $r = \abs{\vect{x}}$, $\uvec{n} = \vect{x}/r$, and write $[f]$ for
a quantity evaluated at the retarded time $t - r/c$. A point dipole is the charge density
$\rho = -\vect{p}(t)\cdot\nabla\delta^3(\vect{x})$ and the current density
$\vect{J} = \dot{\vect{p}}(t)\,\delta^3(\vect{x})$; they conserve charge,
$\partial_t\rho + \nabla\cdot\vect{J} = -\dot{\vect{p}}\cdot\nabla\delta^3 + \dot{\vect{p}}\cdot\nabla\delta^3 = 0$.
In the retarded potential, integrating by parts moves the gradient onto
$\vect{p}(t - \abs{\vect{x}-\vect{x}'}/c)/\abs{\vect{x}-\vect{x}'}$, whose gradient in
$\vect{x}'$ is minus its gradient in $\vect{x}$; so $\phi = -\nabla\cdot([\vect{p}]/r)$. With
$\nabla(1/r) = -\uvec{n}/r^2$ and, for any retarded quantity, $\partial_i[f] = -n_i[\dot f]/c$:
\begin{equation}\label{eq:potentials}
  \phi = \frac{\uvec{n}\cdot[\vect{p}]}{r^2} + \frac{\uvec{n}\cdot[\dot{\vect{p}}]}{c\,r},
  \qquad
  \vect{A} = \frac{[\dot{\vect{p}}]}{c^2 r}.
\end{equation}
Differentiating the same way ($\nabla\times[\vect{f}] = -\uvec{n}\times[\dot{\vect{f}}]/c$,
$\partial_i n_j = (\delta_{ij} - n_in_j)/r$) gives the fields
\begin{align}
  \vect{E} &= \frac{3\uvec{n}(\uvec{n}\cdot[\vect{p}]) - [\vect{p}]}{r^3}
    + \frac{3\uvec{n}(\uvec{n}\cdot[\dot{\vect{p}}]) - [\dot{\vect{p}}]}{c\,r^2}
    + \frac{\uvec{n}\times(\uvec{n}\times[\ddot{\vect{p}}])}{c^2 r}, \label{eq:fields}\\
  \vect{B} &= \frac{1}{c^2}\left(\frac{[\dot{\vect{p}}]\times\uvec{n}}{r^2}
    + \frac{[\ddot{\vect{p}}]\times\uvec{n}}{c\,r}\right), \nonumber
\end{align}
with $\dot{\vect{p}} = -\omega\vect{p}_0\sin\omega t$ and
$\ddot{\vect{p}} = -\omega^2\vect{p}_0\cos\omega t$: the near (static-dipole) term, the
induction term and the radiation term. They are \Jackson{§9.2}'s fields in the time domain and in
the units of \xsecref{foundations}{sec:units}; as $c\to\infty$ only the first remains, the
static dipole field of the instantaneous $\vect{p}(t)$. With $\vect{p}_0$ in the plane $z = 0$,
$\vect{E}$ lies in that plane and $\vect{B}$ along $\uvec{z}$ there, so the ion stays in the
plane.

Two checks of \eqref{eq:fields}, by central differences of step $10^{-4}$ at three points and
times: the fields satisfy Maxwell's equations \xeqref{foundations}{eq:maxwell} in vacuum (all four,
the largest residual $\val{maxwellRes}$ of the largest derivative entering it), and they are the
derivatives of the potentials \eqref{eq:potentials} (to $\val{potentialRes}$), as the
differences' error of about $(10^{-4}/r)^2$ predicts.

\subsection{The flight}\label{sec:flight}

The ion moves by \xeqref{foundations}{eq:eom} with the force
$q(\vect{E} + \vect{v}\times\vect{B})$ of \eqref{eq:fields} at its position and time; in the
plane, $F_x = q(E_x + v_yB_z)$ and $F_y = q(E_y - v_xB_z)$. The field depends on time, so the
energy is not conserved: the field's work changes it,
\begin{equation}\label{eq:work}
  \deriv{}{t}\left(\gamma mc^2\right) = q\,\vect{E}\cdot\vect{v},
  \qquad
  T(t) - T_0 = \int_0^t q\,\vect{E}\cdot\vect{v}\,\dd t',
\end{equation}
(\xsecref{foundations}{sec:energy}: the magnetic force does no work), which the march carries as
a check. The equations have no closed form; they are marched by
\xeqref{foundations}{eq:rk4} at $h = \val{H}$ with two more states, the work \eqref{eq:work} and
the radiated energy $\int P\,\dd t$ of \xeqref{foundations}{eq:larmor} (with
$\vect{a} = (\vect{F} - \vect{v}(\vect{v}\cdot\vect{F})/c^2)/(\gamma m)$). The entry is located on
the cubic Hermite interpolant by bisection and the error estimated against the step $2h$, as in
\xsecref{stray_field}{sec:march}.

\subsection{The kick to first order}\label{sec:kick}

The deflection is small, so the force can be taken along the undeflected line
$x = x_0 + v_0t$, $y = y_0$; to first order in the field the ion gains the transverse momentum
and height
\begin{equation}\label{eq:first-order}
  p_y^{(1)}(t) = \int_0^t q\left(E_y - v_0B_z\right)\!\big(x_0 + v_0t', y_0, t'\big)\,\dd t',
  \qquad
  y^{(1)}(t) = y_0 + \frac{1}{\gamma_0 m}\int_0^t p_y^{(1)}(t')\,\dd t'
\end{equation}
($F_y \perp \vect{v}$, so $\dd v_y/\dd t = F_y/(\gamma m)$ to this order). The field at a point
at distance $b$ from the antenna oscillates as the dipole did a time $b/c$ earlier, and the ion
spends about $b/v_0$ near it. If $b/v_0$ were short against $1/\omega$, the antenna would act as
a static dipole of the moment it has as the ion passes: the kick's sign and size follow
$\cos\omega(t_\mathrm{pass} - b/c)$, and that is the level's point. Here $b/v_0$ is a few
times $1/\omega$, so the oscillation also averages part of the kick away;
\eqref{eq:first-order} includes that. Its error is of second order: the field sampled on the
deflected path instead of the line.

\section{The reference solution by hand}\label{sec:hand}

\begin{procedure}
  \item\label{st:launch} \emph{The field at the launch.} The ion starts $\val{rLaunch}$ above the
    antenna, on its axis ($\uvec{n} = \uvec{y} \parallel \vect{p}_0$). The retarded time is
    $t_\mathrm{r} = -r/c = \val{tRet}$, the phase $\omega t_\mathrm{r} = \val{phaseRet}$, so
    $[p_y] = p_0\cos\omega t_\mathrm{r} = \val{pRet}$ and $[\dot p_y] = \val{pdRet}$. On the axis
    $3\uvec{n}(\uvec{n}\cdot\vect{a}) - \vect{a} = 2\vect{a}$ and
    $\uvec{n}\times(\uvec{n}\times\vect{a}) = 0$: the near term of \eqref{eq:fields} is
    $2[p_y]/r^3 = \val{eNear}$, the induction term $2[\dot p_y]/(cr^2) = \val{eInd}$, the
    radiation term $\val{eRad}$, so $E_y = \val{eYLaunch}$; $\vect{B}$ vanishes on the axis
    ($[\dot{\vect{p}}]\times\uvec{n} = 0$; the orientation's $\cos 90^\circ$, rounded, leaves
    $p_{0x} = \val{ampCross}$). The ion is pushed up with $qE_y = \val{fLaunch}$, accelerating at
    $qE_y/(\gamma_0m) = \val{accLaunch}$.
  \item\label{st:march} \emph{The march.} \S\ref{sec:flight} from the launch, at
    $h = \val{H}$; every 4 time units:

    \smallskip
    \centerline{\val{march}}
    \smallskip

    The ion enters the detector through its \val{faceRef} face at $(\val{xRef}, \val{yRef})$ at
    $t = \val{tRef}$, with $\vect{p} = (\val{pxRef}, \val{pyRef})$ (step estimate
    $\val{errRef}$).
  \item\label{st:work} \emph{The energy.} The field did the work $\val{workRef}$ on the ion
    \eqref{eq:work}, equal to the change of its kinetic energy to $\val{wDevRef}$ along the march.
    The ion radiated $\val{radRef}$ (\ref{as:radiation}).
  \item\label{st:first} \emph{The first-order kick.} \eqref{eq:first-order} by the same
    integrator along the line, to $t^* = (x_\mathrm{det} - x_0)/v_0 = \val{tStar}$ when the line
    reaches the detector's face: $p_y^{(1)} = \val{pyFirst}$ and $y^{(1)} = \val{yFirst}$, which
    differs from the march's $\val{yRef}$ by $\val{firstDiff}$, $\val{firstRel}$ of the lift
    $\val{liftRef}$: the second-order size.
\end{procedure}

\begin{result}[The reference flight]
The ion enters the detector at $(\val{xRef}, \val{yRef})$ at $t = \val{tRef}$, lifted by
$\val{liftRef}$ by the antenna's push near the launch. The level is solved.
\end{result}

\subsection{The energy and the radiation}\label{sec:work}

The field's work is not a one-way transfer: it rises while the antenna pushes the ion near the
launch, falls as the field reverses, and far from the antenna swings with the drive about a small
negative value; at the entry it is $\val{workRef}$ (step~\ref{st:work}).
The radiated energy, $\val{radRef}$, is a fraction $\val{radFrac}$ of $T_0$, far below the
$10^{-10}$ the game allows to neglect (\xsecref{foundations}{sec:radiation}).

\section{Comparison with the game}\label{sec:comparison}

\comparisontable

The trajectories' differences are a few $10^{-12}$ to $\val{largest}$, the game's verified
accuracy (Foundations §6); our step estimates are below $10^{-12}$. The tolerance, $10^{-9}$, is
ten times the larger of the two computations' estimates, or the game's own preview--verification
gap where that is larger: for the flight that leaves the world through the bottom edge, the
game's two runs differ by $\val{gapSame}$ in $x_\mathrm{e}$. The radiated energy agrees to below
$10^{-9}$ relative:
the game integrates the power with three-point Gauss--Legendre on each of its steps (sixth
order), we as a state of the march.

\section{Other solutions}\label{sec:others}

On the empty board the ion flies straight along $y = \val{yZero}$ and leaves the world's right
edge at $t = (x_\mathrm{max} - x_0)/v_0 = \val{bareT}$. A solution must give it an upward kick of
about the reference's, and what decides the kick's sign is the antenna's phase when the ion
passes it (\S\ref{sec:kick}). The table's other placements show it with an antenna along $+y$
two cells below the line:
\begin{itemize}
  \item At $(\val{xAnti}, \val{yAnti})$ the line passes the antenna at $t \approx \val{tPassAnti}$,
    when $\cos\omega(t - b/c) = \val{cosPassAnti}$: in phase ($p_0 = \val{pSame}$) the dipole points
    down then and the ion is pushed down (\val{outSame} at $(\val{xEndSame}, \val{yEndSame})$); in
    the opposite phase ($p_0 = \val{pAnti}$) it is pushed up into the detector, at
    $y = \val{yEndAnti}$.
  \item At $(\val{xLate}, \val{yLate})$ it passes at $t \approx \val{tPassLate}$, with
    $\cos\omega(t - b/c) = \val{cosPassLate}$: in phase ($p_0 = \val{pLate}$) the push is up again, to
    $y = \val{yEndLate}$.
\end{itemize}
Other positions, orientations and amplitudes solve it too, each with the phase that pushes the
ion up as it passes; a map of all solutions would take the game's verdicts, which this document
does not reproduce.

\section{What the game computes}\label{sec:engine}

The game evaluates \eqref{eq:fields} for each antenna, exactly, at the retarded time
(\physref{2.4}), and integrates \xeqref{foundations}{eq:eom} with its adaptive eighth-order
Runge--Kutta method (\physref{3}, \physref{5}, \physref{6}, \physref{7}); the flight is the same
mathematical problem as here, and the differences in the comparison are the two integrations'
errors. It integrates the radiated power as a check and neglects it (\ref{as:radiation}).

\section{Further reading}\label{sec:reading}

The retarded potentials: \Griffiths{Ch.~10}. The oscillating dipole's near, induction and
radiation zones: \Jackson{§9.2}, \Griffiths{Ch.~11}. The flights before this one:
\levelref{stray_field} (a uniform field switched on and off), \levelref{first_bend} (one static
charge).

\levelfileappendix

\end{document}
